\section{Finite schedules for the two splits}\label{schedule:section}

We now turn the three preceding tools into fixed positions at which the
prefix and suffix may read. There will be early cuts that encode metadata
by prefix length, and later cuts that use the tag clock. Every set of cuts
is finite and is fixed using only $M$ and $\Sigma$, before the input word
is read.

The order of selection is part of the construction. We first list all
possible metadata and fix residue classes for their copies. A provisional
clock then realizes the early positions. Those positions determine how
widely markers must be separated. Only after this separation is known do
we choose the prime factors of the final clock. Its approximation is made
on already fixed finite sets, so no earlier length has to be enlarged.
The remaining choices are successive lower bounds on search widths and
gaps. This section verifies both existence and the reading-side properties
of that schedule.

Recall that $V=\mathbf F_2^M$ and $U=V\oplus\mathbf F_2$. The second
space accommodates homogeneous lifts. Positions and cut offsets are
integers; search intervals include both endpoints, while inspected letter
intervals have the half-open convention of Section~\ref{marker:section}.

\subsection{A finite set of metadata and positions}\label{schedule:choices}

Apply Lemma~\ref{alg:shift-table} with entry alphabet $M$ and state spaces
$V$ and $U$, and choose the resulting positive tuple lengths $g$ and $g'$,
respectively.  For the early schedule, a \emph{kind} is a tuple
\begin{equation}\label{schedule:kinds}
 (j,x,y,\mathbf d_{-j},\psi),\qquad
 1\le j\le g',\quad x,y\in U,\quad
 \mathbf d_{-j}\in M^{g'-1},\quad
 \psi:M\longrightarrow\operatorname{End}_{\rm lin}(U).
\end{equation}
The list $\mathbf d_{-j}$ retains the indices and order of all entries
except $j$. The function $\psi$ will record how later letters turn the
monoid value of the early region into a linear update on $U$. There is no
feasibility restriction on any field: in particular, every total function
$\psi$ of the indicated type occurs.
Let $\kappa$ count these kinds. Make $\mu$ distinguishable copies of each
kind, for a positive integer $\mu$ to be fixed below, and call the resulting
roster $\mathcal P$. It has $N=\kappa\mu$ slots. Order them in consecutive
blocks according to $j=1,\ldots,g'$. The copies will let us lose some
positions to false-origin ambiguity without losing any kind altogether.

Set
\begin{equation}\label{schedule:slot-residues}
 B=20(N+1)^3,
 \qquad r_i=4N^2i+i^2\pmod B\quad(1\le i\le N),
\end{equation}
and reserve residue $r_i$ for slot $i$.  The ordered differences $r_i-r_j$,
for $i\ne j$, are all nonzero and distinct modulo $B$.  To check this,
suppose first that $i>j$.  As an integer the difference is
\[
 (i-j)4N^2+(i-j)(i+j).
\]
The second summand lies strictly between $0$ and $2N^2$, so division by
$4N^2$ determines $i-j$.  The remaining summand then determines $i+j$,
and hence $i,j$.  Positive differences are distinct, as are their negatives.
All signed differences lie in an interval of length less than $B$, and no
positive difference equals a negative one.  Reduction modulo $B$ therefore
preserves their distinctness and nonzero values.

The later clock uses a separate tag set $\Lambda$. A main tag has the form
\[
 (j,u,A,C,\mathbf d_{-j},X,Y),\qquad 1\le j\le g,\quad 0\le u<B,
\]
where $A,C\in M$, $\mathbf d_{-j}\in M^{g-1}$, and $X,Y\in V$.
A periodic tag has the form
\[
 (j,u,m_p,X,Y),\qquad 0\le j\le g,\quad 0\le u<B,
\]
where $m_p\in M$ and $X,Y\in V$.
The form itself is part of the tag.  The fields $X,Y$ of a tag $\alpha$ are
denoted $X(\alpha),Y(\alpha)$.  In particular, the clock period and all
marker inspection radii are absent from the tag fields.  If $m=|\Lambda|$,
then
\[
 m=B|V|^2\bigl(g|M|^{g+1}+(g+1)|M|\bigr).
\]
Choose $\mu$ large enough that
\begin{equation}\label{schedule:copies}
 \mu>20\bigl(c_{\rm cl}\log(2m)+4(g+1)\bigr)+10,
\end{equation}
where $c_{\rm cl}$ is the constant in Lemma~\ref{clock:robust}.
Although $m$ depends on $\mu$, this requirement is not circular.
The numbers $\kappa,g,g'$ and the spaces $V,U$ were fixed before $\mu$.
The displayed formulas make $B$ polynomial in $\mu$ and $m$ a constant
multiple of $B$. The right side of \eqref{schedule:copies} thus grows
logarithmically, while the left side grows linearly. Choose and fix one
such $\mu$. From now on the roster, all residue classes, and the entire
tag set are fixed; no clock period or inspection radius is a tag field.

\subsection{The provisional clock and the early schedule}

Apply Lemma~\ref{clock:robust} to $\Lambda$ and the roster.  Write its torus
as $\mathcal T$, its sets as $I_\alpha,J_\alpha,W$, its roster targets as
$v_p$, and its common ball radius as $\rho>0$.  We first select a provisional
additive map $\varphi^0:\mathbf Z\to\mathcal T$.  In each circle coordinate
choose its speed to be $1/n_i^0$, where the $n_i^0$ are distinct sufficiently
large primes greater than $B$, and put
\[
 R=\prod_i n_i^0.
\]
The set
\begin{equation}\label{schedule:provisional-net}
 \{\varphi^0(aB):0\le a<R\}
\end{equation}
is the full product grid with coordinate orders $n_i^0$.  Indeed,
multiplication by $B$ is invertible modulo every $n_i^0$, and the Chinese
remainder theorem prescribes all coordinate residues independently.  The
same assertion holds for any $R$ consecutive representatives of a fixed
residue class modulo $B$.  Choose these primes so large that the grid has
covering radius less than $\rho/3$ in the maximum circular metric, and is
fine enough to approximate each target $v_p$ strictly inside the perturbation
tolerances of Lemma~\ref{clock:robust}.

Put $w_0=BR$.  Realize the ordered roster by positive increasing integers
$p_1,\ldots,p_N$ such that
\begin{equation}\label{schedule:stencil-properties}
 p_i\equiv r_i\pmod B,
 \qquad \varphi^0(p_i)\text{ is within the chosen tolerance of }v_{p_i},
\end{equation}
and all ordered nonzero differences $p_i-p_j$ are at distance greater than
$2w_0+2$ from each other and from $0$.  These requirements can be met
inductively.  For each slot, the grid property supplies an admissible
integer in its reserved class, and adding multiples of $BR$ leaves both
that class and its provisional clock value unchanged.  At the next
insertion, place the new point so far to the right that all new positive
differences exceed every old positive difference by more than $2w_0+2$.
The differences between two of these new differences are differences of
previous points, which already have the required separation.  Their
negatives satisfy the same conditions.  Also leave an integer gap between
successive blocks of slots.

From this point, $\mathcal P$ denotes this set of integer offsets, with each
offset still carrying its original kind and copy label.  Write
$h=\operatorname{diam}\mathcal P$.  Choose fixed integer boundary offsets
\[
 0=\beta_0<\beta_1<\cdots<\beta_{g'}
\]
so that every slot of block $j$ lies strictly between $\beta_{j-1}$ and
$\beta_j$.  For a word beginning at $s$, the corresponding early boundaries
are
\begin{equation}\label{schedule:stencil-boundaries}
 b'_j=s+\beta_j\qquad(0\le j\le g'),
\end{equation}
and its early candidate cuts are $s+p$, $p\in\mathcal P$.

\subsection{The final clock and inner marker searches}

The early positions are now fixed.  We next choose the inner marker rule,
and only then the prime factors of the final clock.  Choose integers
\begin{equation}\label{schedule:inner-marker-parameters}
 d>10(2h+2w_0+5),\qquad K>3d^2,
\end{equation}
and obtain an inner marker rule of radius $r\ge K$ from
Lemma~\ref{marker:local}.  Put $N_M=|M|!$.

In each coordinate, approximate the provisional speed $1/n_i^0$ by
$a_i/n_i$, where the $n_i$ are distinct primes greater than $d,N_M,B$ and
$1\le a_i<n_i$.  Arbitrarily accurate such approximations exist by taking
the primes sufficiently large and rounding $n_i/n_i^0$ to an integer.
Let $\varphi$ be the resulting additive map and put
\begin{equation}\label{schedule:final-period}
 n=\prod_i n_i.
\end{equation}
The map $\varphi$ factors through $\mathbf Z/n\mathbf Z$.  Approximate
closely enough on the already fixed finite sets
$\mathcal P$ and $\{aB:0\le a<R\}$ that
\begin{equation}\label{schedule:clock-roster}
\begin{aligned}
 &\#\{(p,p')\in\mathcal P^2:p\ne p',\
 a+\varphi(p')-\varphi(p)\in W\}\\
 &\qquad\le c_{\rm cl}\log(2m)\qquad(a\in\mathcal T).
\end{aligned}
\end{equation}
\begin{equation}\label{schedule:final-net}
 \{\varphi(aB):0\le a<R\}
 \text{ has covering radius less than }\rho.
\end{equation}
For the first assertion, keep every $\varphi(p)$ within the clock's
single allowed perturbation tolerance of its target. For the second,
keep every value at an argument $aB$, $0\le a<R$, within the unused
margin between the provisional covering radius and $\rho$. Only finitely
many arguments are involved, so one sufficiently accurate approximation
satisfies both requirements. Every translate of this fixed net meets
every clock ball of radius $\rho$.

The construction now has the two properties it will use separately:
the clock period has no small prime factors, while the useful $R$-term
net and $w_0=BR$ keep their earlier sizes. Replacing $R$ by the final
period here would lose this independence and is unnecessary.

Choose positive integers $L_{\rm lag}$ and $H_{\rm ext}$ with
\begin{equation}\label{schedule:lag}
 B\mid L_{\rm lag},\qquad
 L_{\rm lag}>r+(|M|+2)d,\qquad
 H_{\rm ext}>L_{\rm lag}+ndN_MB+2d.
\end{equation}
Choose positive widths $H_u$, indexed by $0\le u<B$, each divisible by
$B$ and greater than $2d$, with pairwise distances
\begin{equation}\label{schedule:width-spacing}
 |H_u-H_v|>2h+2\max(H_{\rm ext},w_0)+2
 \qquad(u\ne v).
\end{equation}
For each fixed anchor in the false-origin count, distinct displacement
residues will give distinct width indices. This spacing will separate the
right endpoints of searches whose left endpoints lie in one short interval.
Finally choose a positive multiple $G$ of $2B$ such that
\begin{equation}\label{schedule:block-spacing}
 G>10\bigl(\max_u H_u+H_{\rm ext}+r+K+d+w_0+1\bigr).
\end{equation}
All these requirements are lower bounds on successively chosen integers,
so they are simultaneously satisfiable.

An episode beginning at $s$ will have a sampling origin $t=s+L$, where the
fixed positive lag $L$ will be chosen after all bounded windows are known.
For any integer origin $t$, define
\begin{equation}\label{schedule:inner-searches}
 o_j=(j+1)G,\qquad z_j(t)=t+o_j,\qquad
 f_j^u(t)=\text{the first inner marker in }[z_j(t),z_j(t)+H_u],
\end{equation}
for $0\le j\le g$ and $0\le u<B$.  The value is $\bot$ when the search
finds no marker.  Equality of search values means equality of absolute
positions, or that both values are $\bot$.

\begin{lemma}[Inner cut schedules]\label{schedule:inner-cuts}
For $1\le j\le g$, put $c_j=(j+\tfrac12)G$.  The main cuts of block $j$
are
\begin{equation}\label{schedule:main-cuts}
 k=t+c_j+aB,\qquad 0\le a<R.
\end{equation}
At any such cut, all inner searches through index $j-1$, including their
marker inspections, lie before $k$.  The searches from index $j$ onward
lie after $k$ even when performed at any of the hypothesized origins
$k-c_j-a'B$, $0\le a'<R$, and with any width index $u$.

For $0\le j\le g$ and $0\le u<B$, the periodic cuts are
\begin{equation}\label{schedule:periodic-cut-formula}
 k=z_j(t)+H_u+L_{\rm lag}+aB,
 \qquad 0\le a<ndN_M.
\end{equation}
The whole search for $f_j^u(t)$, with its inspections, lies before every
one of these cuts.  Every offset $k-t$ in either schedule is divisible
by $B$.
\end{lemma}

\begin{proof}
The rightmost inspection for a search through $j-1$ is at most
$t+jG+\max_uH_u+r$, strictly before $t+(j+\tfrac12)G$ by
\eqref{schedule:block-spacing}.  For a search from $j$ onward at
$\tau=k-c_j-a'B$, its leftmost possible inspection is at least
\[
 \tau+(j+1)G-r=k+\tfrac12G-a'B-r>k,
\]
because $a'B<w_0$ and $G/2>w_0+r$.  The periodic assertion follows from
$L_{\rm lag}>r$.  Divisibility follows from $2B\mid G$ and the choices of
$H_u,L_{\rm lag}$.
\end{proof}

Let $\mathcal H$ be the full finite set of offsets $k-t$ in
\eqref{schedule:main-cuts} and \eqref{schedule:periodic-cut-formula}, over
all their indicated indices.  Choose an integer
\begin{equation}\label{schedule:end-stability-width}
 w>\max(2h,\operatorname{diam}\mathcal H)+1.
\end{equation}
The next observation records exactly how the periodic schedule will
supply clock choices when an inner marker search is absent.

\begin{lemma}[Periodic cuts with a fixed prefix product]
\label{schedule:periodic-cuts}
Fix $j,u$ and an origin $t$.  Suppose the interval
$[z,z+H_u+H_{\rm ext}]$, where $z=z_j(t)$, contains no inner marker and
all required inspections are available.  There is a period $d'<d$ such
that, with $k_0=z+H_u+L_{\rm lag}$, the cuts
\begin{equation}\label{schedule:periodic-progression}
 k=k_0+a d'N_MB,\qquad 0\le a<n,
\end{equation}
belong to \eqref{schedule:periodic-cut-formula}.  For a fixed beginning
$s\le z$, $T_{s,k}$ is constant along this progression.
Moreover $\varphi(k-t)$, as $a$ varies, traverses a translate of the full
image of $\varphi$ and therefore includes a translate of the net in
\eqref{schedule:final-net}.
\end{lemma}

\begin{proof}
Lemma~\ref{marker:periodic} gives one periodic continuation, of some
period $d'<d$, throughout
\[
 [z+d,z+H_u+H_{\rm ext}-d+K).
\]
The $|M|$ periods immediately before $k_0$ lie in this interval:
$k_0-|M|d'>z+d$ follows from \eqref{schedule:lag}.  All the cuts in
\eqref{schedule:periodic-progression} lie there as well, because
\[
 L_{\rm lag}+(n-1)d'N_MB
 <L_{\rm lag}+ndN_MB<H_{\rm ext}-2d.
\]
They belong to the stated periodic schedule since $a d'N_M<ndN_M$.

Let $c\in M$ be the product of one period in phase at $k_0$.  Among
$1,c,\ldots,c^{|M|}$ two powers agree.  If their exponent difference is
$v$, then $1\le v\le |M|$, $v\mid N_M$, and multiplication by further
powers gives
\[
 c^{i+N_M}=c^i\qquad(i\ge |M|).
\]
It follows that every cut in \eqref{schedule:periodic-progression} has
\[
 T_{s,k}
 =T_{s,k_0-|M|d'}c^{|M|+aN_MB}
 =T_{s,k_0-|M|d'}c^{|M|}.
\]
Finally, each prime dividing $n$ is greater than $d,N_M,B$.  Thus
$\gcd(d'N_MB,n)=1$, so the increments $a d'N_MB$, for $0\le a<n$, run
through all residues modulo $n$.  Apply the additive map $\varphi$.
\end{proof}

The schedule is now complete except for its placement inside an episode.
Main cuts separate the required marker products, and periodic cuts provide
a fixed prefix product while sweeping the clock image. Both schedules
have offsets divisible by $B$. The next section puts all their windows,
including the displaced windows used by false suffix hypotheses, strictly
inside each episode. A separate marker rule there determines the end.
